% General comparison lemma; numerical q7 conclusion is a separate gate.
\subsection{Small projective certificates for act-cone squares}

\begin{lemma}[The literal act-cone kernel]
\label{C:act-cone-literal-bar-kernel}
Let $M$ be a monoid, $S$ a nonempty right $M$-set, $k$ a field,
$R=kM$, and $L=\ker(kS\to k)$.
Write $C(M,S)$ for the two-object act category.
The kernel of the normalized-chain map
\[
 C_*(BC(M,S);k)\longrightarrow C_*(BC(M,*);k)
\]
is naturally the suspension of the normalized two-sided bar
$B_*(L,R,k)$.  In degree $n+1$ the comparison sends
\[
 w[a_1|\cdots|a_n]\longmapsto
 (-1)^{n(n-1)/2}\sum_{s\in S}w_s
       [a_n|\cdots|a_1|s],\qquad w=\sum_s w_s s.
\]
Consequently it induces
$\operatorname{Tor}_n^R(L,k)\cong\widetilde H_{n+1}(BC(M,S);k)$.
\end{lemma}
\begin{proof}
The two categories have the same objects and endomorphisms.
A nondegenerate simplex which changes objects consists of monoid
arrows followed by exactly one arrow labelled by a state.
The kernel on its state coefficients is precisely $L$.
All other normalized chains map identically.  This identifies the
kernel degreewise with the displayed two-sided bar.

Let $T_n$ denote the displayed map without its sign.
The last category face cancels because $\sum_s w_s=0$.
The face next to the state is its right $M$-action; the remaining
faces multiply consecutive monoid letters or remove the other end.
Reversing the letters gives
$\partial T_n=(-1)^nT_{n-1}d$.
The sign in the statement changes this to the usual suspended
differential $-d$.  The chain map is therefore a literal isomorphism
onto the kernel.  The quotient category $C(M,*)$ has a terminal
object and contractible classifying space.  The exact sequence of
normalized chains gives the asserted homology isomorphism.
\end{proof}

\begin{proposition}[A projective comparison for the affine square]
\label{C:act-projective-affine-square}
Suppose $M$ has minimal band $I\times J$ acting constantly on $S$.
Choose base elements $i_0\in I$ and $s_0\in S$ as in the
constant-action compression, and put
$V=\widetilde{kI}$ and $W=\widetilde{kS}=L$.
Let $F:W\times V\to k$ be a balanced bilinear form and let $u_F$
be its affine degree-two class on $C(M,S)$.
Set
\[
 \alpha_a=e_{\sigma_a(i_0)}-e_{i_0},\qquad
 \beta_b=e_{s_0b}-e_{s_0}.
\]
On the two-sided bar of Lemma~\ref{C:act-cone-literal-bar-kernel},
the pullback of $u_F^2$ in degree three is, up to the one common
minus sign supplied by that lemma,
\[
 w[a|b|c]\longmapsto F(w,\alpha_a)F(\beta_b,\alpha_c).
\]

Let $P_\bullet\to k$ be a projective resolution of the trivial
left $R$-module, with each displayed summand a literal ideal
$R\varepsilon$, $\varepsilon^2=\varepsilon$.
Projective terms $P_0,\ldots,P_3$ with exactness through $P_2$
(so $P_3$ surjects onto $\ker d_2$), together with their comparison
to the normalized free left bar resolution, suffice to certify that
this square is a coboundary by exhibiting a primitive on
$L\otimes_RP_\bullet$.
No termination of the undisplayed resolution is required.
\end{proposition}
\begin{proof}
The affine representative on an old pair $g$ then $f$ is
$F(\beta_f,\alpha_g)$; on an old arrow $g$ followed by the state
$s$ it is $F(e_s-e_{s_0},\alpha_g)$.
These are the normalized cocycles supplied by the actual
constant-action compression of
Theorem~\ref{C:constant-minimal-act-compression}.
The two Alexander--Whitney halves of
$[c|b|a|s]$ give
\[
 F(\beta_b,\alpha_c)F(e_s-e_{s_0},\alpha_a).
\]
Summing with coefficients $w_s$ gives the formula because their
sum is zero.  In degree three the sign of the preceding lemma is
$-1$; multiplying a cochain and its primitive by $-1$ does not
change whether the class is zero.

Here is a completely explicit comparison on a projective prefix.
The free left bar resolution has basis $(r;a_1,\ldots,a_n)$,
where each $a_i\ne1$, and its underlying augmented complex has
contraction
\[
 s(r;a_1,\ldots,a_n)=
 \begin{cases}(1;r,a_1,\ldots,a_n),&r\ne1,\\0,&r=1.
 \end{cases}
\]
The degree-minus-one contraction sends $1$ to $(1;\varnothing)$.
For each idempotent generator define recursively
\[
 \phi_n(\varepsilon)
   =\varepsilon\,s\bigl(\phi_{n-1}(d\varepsilon)\bigr),
\]
using the augmentation in degree zero.
The argument of $s$ is a cycle.  It is fixed by $\varepsilon$,
since $d\varepsilon$ is fixed by that idempotent and the previous
comparison is $R$-linear.  Thus
$d\phi_n(\varepsilon)=\phi_{n-1}(d\varepsilon)$.
The identity
$\varepsilon\phi_n(\varepsilon)=\phi_n(\varepsilon)$ makes the
extension $r\varepsilon\mapsto r\phi_n(\varepsilon)$ well defined
and $R$-linear.  This proves all comparison and module equations
on a generating set, and hence everywhere.

After tensoring, a term $(r;a,b,c)$ in
$\phi_3(\varepsilon)$ contributes
\[
 F(wr,\alpha_a)F(\beta_b,\alpha_c).
\]
This gives the square cochain on the small tensor complex.
If its row vector is the product of a row vector in degree two
with the tensor differential $d_3$, that vector is an explicit
degree-three primitive for the square on the act cone.
The projective prefix extends to a full projective resolution by
first covering $\ker d_3$ and then continuing to cover kernels;
the same contraction extends
$\phi$.  The comparison of these resolutions computes the derived
tensor product in every degree.  Therefore the certified primitive
has the asserted meaning in ordinary cohomology, independently of
whether the remaining resolution is finite.

For clarity, tensoring a map
$R\varepsilon\to R\eta$ whose generator has value $z$ uses
$w\mapsto wz$.  Its literal support is
$z=\varepsilon z=z\eta$.
On column-coordinate matrices for a right module this says
$R_z=E_\eta R_z E_\varepsilon$.
The genuine tensor terms are the images $W\varepsilon$.
Extending the matrix to ambient copies of $W$ preserves its rank,
and a supported cochain is a coboundary exactly when it belongs
to the ambient differential's row space.  This is the finite
linear-algebra test used by the certificates.
\end{proof}

\begin{remark}[What a finite-field square certificate proves]
\label{C:act-projective-square-certificate-scope}
In the application over $\mathbb F_{65521}$, the exact contracted
projective prefix gives
$\dim\operatorname{Tor}_0^{kM/kK}(W,V)=1$.
A nonzero balanced form is therefore the nonzero degree-two
cohomology class, by the affine representative and compression
comparison.  A primitive from the preceding proposition proves
that this actual field class has square zero and excludes
$\mathbb CP^2$.  It does not infer rational rank from a modular
rank, nor infer vanishing of higher homology from a finite prefix.
The routines are
\path{code/a_monoid_bar_prefix_modp.py} and
\path{code/a_act_bar_square_modp.py}; all projective ideals, full
embedded kernels, idempotent supports and sparse comparison chain
equations are retained explicitly.
\end{remark}
